\documentclass[../main.tex]{subfiles}

\begin{document}

\section{Cambios de ideas}

Aumente la cantidad de magnitudes de los vectores de onda para tener una mejor resolución en la gráfica del factor de estructura.
Aún no he hecho la prueba con una estructura cúbica y analizado que mi implementación de un factor de estructura sea correcta.

El principal cambio entre \(t_{o}\) y \(t_{f}\) en el factor de estructura es el aumento de la magnitud de \(S(\vec{q})\) en los mismos valores de \(\vec{q}\).
Entonces me puse a buscar, sin haber comprobado con una estructura cúbica, otras técnicas de análisis para la red, por si el factor de estructura no es un buen observable.

\subsection{Opciones según DeepSeek}

\paragraph{Teoría de percolación}
Relacionar la fracción de enlaces y la conectividad con parámetros de interés.

\paragraph{Ecuación de Smoluchowski con kernel dependiente del potencial.}
Modela la evoluación de la distribución de tamaños agregados\footnote{Puede que funcione para fracciones de empaquetamiento bajos o bajas densidades de crosslinkers.}.

\paragraph{Teoría de perturbación termodinámica de Wertheim para fluidos asociantes.}
Para interacciones con dirección y de corto alcance, formando enlaces físicos transistorios.

\paragraph{Modelo de esferas adhesivas (Baxer) y solución analítica de Percus-Yevick}
Al considerar que el potencial atractivo es de corto alcance, el potencial exponencial se puede mapear a un potencial adhesivo con un parámetro de pegajosidad.

\paragraph{Teoría de acoplamiento de modos para geles atractivos.}
Predice la transición dinámica a un estado no ergódico en función de la fuierza y el alcance de la atracción.

\paragraph{Modelo de energía de paisaje y cinéctica de reacción entre estados.}
Considerar los enlaces como activos o inactivos y formula una ecuación maestra para la red.

\subsection{Notas al aire}

No estoy muy seguro que se está caracterizando.
Si el sistema percola?
La elasticidad?
LA geometría?
No me queda muy claro cuál es la tirada.

Me acuerdo que tiene que ver con impresión de hidrogeles.
Una impresión por extrucción.
Analizar las propiedades mecánicas.

\section{Sugerencias de DeepSeek}

Al no estar seguro que analizar, usé DeepSeek para crear una lista de los observables más comunes que se usan para caracterizar geles, de forma experimental o por simulación computacional.

\subsection{Conectividad y estructura de la red}

\paragraph{Probabilidad de percolación \(P_{\mathrm{perc}}\)}
Fracción de configuraciones que representa un cluster que conecta dos bordes opuestos del sistema.

\paragraph{Fracción de partículas en el cluster que percola.}
Masa del cluster que percola divida entre el número total de partículas.
Cero en la fase sol y finita en el gel.

\paragraph{Distribución de tamaños de clusters}
Número de clusters de tamaño \(s\).
En el punto de gelación muestra una cola de ley de potencias.

\paragraph{Tamaño medio de cluster.}
Segundo momento de la distribución de tamaño de cluster, axcluyendo el cluster percolante.

\paragraph{Longitud de correlación de conectividad.}
Distancia característica de la malla de la red incipiente.
Se extrae del factor de estructura \(S(q)\) o por conteo directo de pares conectados.
Diverge en el punto crítico.


\subsection{Estructura espacial y orden local}

\paragraph{Factor de estructura estático.}
En geles se observa un pico a bajo \(q\) asociado a la longitud de la malla y un pico a \(q\) alto por la escala de la partícula.
La intensidad a bajo \(q\) aumenta drásticamente al gelificar.

\paragraph{Función de distribución radial.}
Probabilidad de encontrar una partícula a distancia \(r\) de otra.
Permite identificar enlaces formados y la extensión de la capa de coordinación, así como cambios en la estructura local al variar la atracción.

\paragraph{Número de coordinación medio \(Z\).}
Promedio de vecinos enlazados por partícula.

\paragraph{Parámetro de orden orientacional}

\subsection{Propiedades dinámicas}

\paragraph{Desplazamiento cuadrático medio.}
La altura de la meseta y la difusividad de largo plazo están ligadas a la rigidez de la red.

\paragraph{Función de correlación de densidad intermedia.}
Decaimiento de las fluctuaciones de densidad.
La congelación del gel se manifiesta como un parámetro de no ergodicidad.
Cuanta estructura queda congelada.

\paragraph{Coeficiente de difusión colectiva y depdencia con el vector de onda.}
Obtenido de \(S(q,t)=S(q)e^{D_{c}q^2 t}\) a bajo \(q\).
Al gelificar, \(D_c\) cae varios órdenes de magnitud.

\subsection{Propiedades mecánicas y reológicas.}

\paragraph{Modulos de almacenamiento y pérdida.}

\paragraph{Viscosidad de corte cero.}

\subsection{Observables termodinámicos y energéticos}

\paragraph{Energía potencial media por partícula.}

\paragraph{Capacidad calorífica configuracional o fluctiaciones de energía.}

\subsection{Observables específicos de percolación}

\paragraph{Distribución de fuerzas en los enlaces.}

\paragraph{Conectividad de loops o circuitos cerrados.}
Número de ciclos independientes.
Permite distinguir entre agregados ramificados y geles con ciclos.

\subsection{Parámetros directos de DeepSeek}

\begin{itemize}
    \item Fracción de partículas en el cluster percolante
    \item Tamaño medio de cluster 
    \item Distribución de tamaños de clusters 
    \item Factor de estructura
    \item Desplazamiento cuadrático medio
    \item Vida media de enlaces
\end{itemize}

Según deepseek, esos parámetros relacionan los parámetros de profundidad del potencial y distancia de interacción con la estructura y estabilidad del gel.

\end{document}
